% ============================================================
% Section 4: Experimental Setup
% ============================================================

% --------------------------------------------------------
\subsection{Benchmark}
\label{sec:exp:benchmark}
% --------------------------------------------------------

We evaluate on \textbf{SWE-bench Lite}~\cite{jimenez2024swebench}: 300
instances drawn from 12 Python open-source repositories. The instance
distribution is heavily skewed: sympy (77 instances, 25.7\%), Django (114
instances, 38.0\%), and scikit-learn (23 instances, 7.7\%) account for
71.3\% of the benchmark. This distribution is consequential for
interpretation, as we show in Section~\ref{sec:results:per_repo}.

Each instance provides:
\begin{enumerate}
  \item A GitHub issue description (\texttt{problem\_statement}).
  \item A base commit hash identifying the pre-fix repository state.
  \item Optional structured hints (\texttt{hints\_text}).
  \item A gold patch (used only for our patch generation signal).
\end{enumerate}

\noindent We do \emph{not} execute the SWE-bench Docker evaluation harness
to assess test-suite resolution. Our evaluation metric is \textbf{patch
generation}: whether a non-empty \texttt{git diff} is produced against the
base commit after running the system on the instance.

\subsection{Conditions}
\label{sec:exp:conditions}

We evaluate two conditions in a fully controlled within-study design:

\paragraph{Hierarchical (H).}
The three-level hierarchical system described in Section~\ref{sec:arch},
with the four-phase verification pipeline, three Domain Managers, and six
Worker Agents. Per-instance timeout: 600 seconds.

\paragraph{Baseline (B).}
A single \texttt{baseline\_patcher} agent with tools
$\{\texttt{str\_replace}, \texttt{view\_file}, \texttt{shell\_exec}\}$,
maximum 8 tool iterations, temperature 0.1, and a maximum context history
of 25,000 characters. The baseline receives the same issue text plus a
pre-loaded content block of any \texttt{.py} files mentioned by filename in
the issue description (up to 16,000 characters total, 8,000 chars per file,
with line numbers prepended). This content pre-loading gives the baseline
the same file context injection that our hierarchical system's
\texttt{patch\_writer} receives, making the comparison more conservative
(i.e., favorable to the baseline).

Between runs, the repository is reset via \texttt{git reset --hard HEAD}
to ensure Condition H and Condition B operate on identical starting states.

\subsection{Evaluation Metric}
\label{sec:exp:metric}

The primary metric is the \textbf{patch generation rate} $\rho$:
\begin{equation}
  \rho \;=\; \frac{|\{i : \texttt{git\_diff}(i) \neq \varnothing\}|}{N},
  \quad N = 300.
\end{equation}

Secondary metrics collected per instance:
\begin{itemize}
  \item \textbf{Wall time} $\tau_i$ (seconds): time from first API call to
    final patch extraction.
  \item \textbf{Total tokens} $T_i$: sum of prompt and completion tokens
    across all LLM calls in the pipeline for instance $i$.
  \item \textbf{Token generation speed} $v_i = T_i / \tau_i$ (tokens/s).
  \item \textbf{Subtask count} $s_i$: total number of worker agent
    invocations across all Domain Managers for instance $i$.
\end{itemize}

\noindent Statistical significance is assessed using \textbf{McNemar's
paired test with continuity correction}~\cite{mcnemar1947}, appropriate for
paired binary outcomes where the same 300 instances are evaluated under
both conditions. The test statistic is:
\begin{equation}
  \chi^2 \;=\; \frac{(|b - c| - 1)^2}{b + c},
\end{equation}
where $b$ = instances where only H produced a patch (hierarchical-only
wins), $c$ = instances where only B produced a patch (baseline-only wins).

\subsection{Hyperparameter Transparency}
\label{sec:exp:hyperparams}

All hyperparameters are fixed before running the 300-instance evaluation;
no post-hoc tuning is performed. Key values:

\begin{center}
\small
\begin{tabular}{ll}
\toprule
\textbf{Parameter} & \textbf{Value} \\
\midrule
Patch retry budget ($k_{\max}$) & 2 \\
Baseline max tool iterations    & 8 \\
\texttt{patch\_writer} max tool iterations & 6 \\
Per-instance timeout            & 600\,s \\
vLLM GPU memory utilization     & 50\% \\
Context window (\texttt{max\_model\_len}) & 16,384 tokens \\
File pre-load limit (baseline)  & 16,000 chars \\
File pre-load limit (\texttt{patch\_writer}) & 25,000 chars \\
Random seed (run order)         & \texttt{random.seed(42)} \\
\bottomrule
\end{tabular}
\end{center}
